\input{header.tex}
\input{newcommands.tex}

\title{Sprechblasen mit Ti\textit{k}Z in \hologo{LaTeX}}
\author{Handreichung (Ergänzung) für MINT-Lehrkräfte}
\date{21. Dezember 2025}
\pagestyle{plain}


% \lehrerblase und \schuelerblase stehen in der universalen newcommands.tex

\begin{document}
\maketitle

Mit der Ti\textit{k}Z-Bibliothek \texttt{shapes.callouts} lassen sich klassische Comic-Sprechblasen (engl. \emph{callouts}) erzeugen: eine (meist abgerundete) Form mit einem \emph{Pointer} bzw. \emph{Tail} (deutsch z.\,B. \emph{Zeiger} oder \emph{Sprechblasen-Zipfel}), der auf eine Stelle zeigt.

\vspace{0.3cm}
\noindent\textbf{Ziel dieser Handreichung:} Sie erstellen Schritt für Schritt Sprechblasen im Stil „abgerundetes Rechteck“ mit Zipfel links (Lehrer) bzw. rechts (Schüler) und lernen, den Zipfel gezielt zu positionieren.

%%%%%  KAPITEL 1  %%%%
\section{Bibliothek einbinden und Grundidee}

Für Callouts benötigen Sie Ti\textit{k}Z und die Bibliothek \texttt{shapes.callouts}:

\begin{lstlisting}[style=latex]
\usepackage{tikz}
\usetikzlibrary{shapes.callouts}
\end{lstlisting}

Eine Sprechblase ist in Ti\textit{k}Z typischerweise ein \texttt{node} mit einer Callout-Form, z.\,B. \texttt{rectangle callout} oder \texttt{ellipse callout}. Der Zipfel wird über den Schlüssel \texttt{callout relative pointer=\{(dx,dy)\}} gesetzt.

\begin{lstlisting}[style=latex]
\begin{tikzpicture}
  \node[rectangle callout, draw, callout relative pointer={(8mm,-6mm)}] {Text};
\end{tikzpicture}
\end{lstlisting}

\begin{ausgabe}
\begin{center}
\begin{tikzpicture}
  \node[rectangle callout, draw] {Text};
\end{tikzpicture}
\end{center}
\end{ausgabe}

\noindent\textbf{Merksatz:} \texttt{callout relative pointer} verschiebt den Zipfel relativ zur \emph{Mitte} der Blase: \\
\(\;\;\)positives \(x\): nach rechts, negatives \(x\): nach links; positives \(y\): nach oben, negatives \(y\): nach unten.

%%%%%  KAPITEL 2  %%%%
\section{„Echte“ Comic-Optik: abgerundetes Rechteck}

Für eine typische Sprechblase im Unterrichtsmaterial ist \texttt{rectangle callout} in Kombination mit \texttt{rounded corners} sehr praktikabel:

\begin{lstlisting}[style=latex]
\begin{tikzpicture}
  \node[
    rectangle callout,
    rounded corners=6pt,
    draw,
    thick,
    fill=blue!6,
    align=left,
    inner sep=3.5mm,
    callout relative pointer={(-12mm,-7mm)}
  ] {Lehrer: Guten Morgen!};
\end{tikzpicture}
\end{lstlisting}

\begin{ausgabe}
\begin{center}
\begin{tikzpicture}
  \node[rectangle callout,rounded corners=6pt,draw,thick,fill=blue!6,align=left,inner sep=3.5mm,callout relative pointer={(-12mm,-7mm)}] {Lehrer: Guten Morgen!};
\end{tikzpicture}
\end{center}
\end{ausgabe}

\noindent Hier erzeugt \texttt{rounded corners=6pt} die weichen Ecken, \texttt{fill=blue!6} eine leichte Füllung und \texttt{inner sep} den Innenabstand.

%%%%%  KAPITEL 3  %%%%
\section{Zipfel gezielt setzen: links vs. rechts}

In der Praxis möchten Sie den Zipfel oft \emph{spiegeln}: Lehrer links, Schüler rechts. Dazu ändern Sie nur das Vorzeichen der \(x\)-Komponente.

\subsection*{Lehrerblase (Zipfel links)}
\begin{lstlisting}[style=latex]
\begin{tikzpicture}
  \node[rectangle callout, rounded corners=6pt, draw, thick,
        fill=blue!6, inner sep=3.5mm, align=left,
        callout relative pointer={(-13mm,-7mm)},
        text width=0.7\linewidth]
  {\textbf{Einstieg:} Wer kennt Geheimschriften aus der Kindheit?};
\end{tikzpicture}
\end{lstlisting}

\begin{ausgabe}
\begin{center}
\begin{tikzpicture}
  \node[rectangle callout, rounded corners=6pt, draw, thick,
        fill=blue!6, inner sep=3.5mm, align=left,
        callout relative pointer={(-13mm,-7mm)},
        text width=0.7\linewidth]
  {\textbf{Einstieg:} Wer kennt Geheimschriften aus der Kindheit?};
\end{tikzpicture}
\end{center}
\end{ausgabe}

\subsection*{Schülerblase (Zipfel rechts)}
\begin{lstlisting}[style=latex]
\begin{tikzpicture}
  \node[rectangle callout, rounded corners=6pt, draw, thick,
        fill=green!6, inner sep=3.5mm, align=left,
        callout relative pointer={(13mm,-7mm)},
        text width=0.55\linewidth]
  {ROT13, Zahlen statt Buchstaben, \dots};
\end{tikzpicture}
\end{lstlisting}

\begin{ausgabe}
\begin{center}
\begin{tikzpicture}
  \node[rectangle callout, rounded corners=6pt, draw, thick,
        fill=green!6, inner sep=3.5mm, align=left,
        callout relative pointer={(13mm,-7mm)},
        text width=0.55\linewidth]
  {ROT13, Zahlen statt Buchstaben, \dots};
\end{tikzpicture}
\end{center}
\end{ausgabe}

\noindent\textbf{Praxis-Tipp:} Wenn der Zipfel „zu steil“ wirkt, verringern Sie \(|dy|\) (z.\,B. \(-5mm\)). Wenn er zu kurz ist, erhöhen Sie \(|dx|\) (z.\,B. \(16mm\)).

%%%%%  KAPITEL 4  %%%%
\section{Relative vs. absolute Pointer-Koordinaten}

Ti\textit{k}Z kennt zwei gebräuchliche Methoden, den Zipfel zu setzen:

\subsection*{1) Relativ (empfohlen fürs Layout)}
\begin{lstlisting}[style=latex]
callout relative pointer={(dx,dy)}
\end{lstlisting}
Diese Variante ist robust, weil sie unabhängig davon ist, wo die Blase im Dokument steht.

\subsection*{2) Absolut (zeigt auf eine konkrete Stelle)}
\begin{lstlisting}[style=latex]
callout absolute pointer={(x,y)}
\end{lstlisting}

Im nächsten Beispiel zeigt die Blase auf einen Punkt (z.\,B. ein Objekt im Bild):

\begin{lstlisting}[style=latex]
\begin{tikzpicture}
  \fill (0,0) circle (2pt); % Zielpunkt
  \node[rectangle callout, rounded corners=6pt, draw,
        callout absolute pointer={(0,0)},
        fill=yellow!20] at (3,1) {Das ist der Punkt.};
\end{tikzpicture}
\end{lstlisting}

\begin{ausgabe}
\begin{center}
\begin{tikzpicture}
  \fill (0,0) circle (2pt); % Zielpunkt
  \node[rectangle callout, rounded corners=6pt, draw,
        callout absolute pointer={(0,0)},
        fill=yellow!20] at (3,1) {Das ist der Punkt.};
\end{tikzpicture}
\end{center}
\end{ausgabe}

%%%%%  KAPITEL 5  %%%%
\section{Feintuning: Breite, Rundung und Kürzen des Zipfels}

Die wichtigsten Stellschrauben für die Optik sind:

\begin{itemize}
  \item \texttt{callout pointer width=<Länge>}: Breite des Zipfels an der Ansatzstelle
  \item \texttt{callout pointer arc=<Zahl>}: Rundung der Zipfel-Kanten (comic-iger)
  \item \texttt{callout pointer shorten=<Länge>}: kürzt den Zipfel (z.\,B. wenn er zu lang wirkt)
\end{itemize}

\begin{lstlisting}[style=latex]
\begin{tikzpicture}
  \node[rectangle callout, rounded corners=8pt, draw, thick,
        fill=orange!15,
        callout relative pointer={(14mm,-6mm)},
        callout pointer width=10mm,
        callout pointer arc=18,
        callout pointer shorten=1mm]
  {Weicher Zipfel.};
\end{tikzpicture}
\end{lstlisting}

\begin{ausgabe}
\begin{center}
\begin{tikzpicture}
  \node[rectangle callout, rounded corners=8pt, draw, thick,
        fill=orange!15,
        callout relative pointer={(14mm,-6mm)},
        callout pointer width=10mm,
        callout pointer arc=18,
        callout pointer shorten=1mm]
  {Weicher Zipfel.};
\end{tikzpicture}
\end{center}
\end{ausgabe}

%%%%%  KAPITEL 6  %%%%
\section{Komfort: Lehrer-/Schüler-Befehle zum Wiederverwenden}

Damit Sie in Ihren Stundenplanungen schnell arbeiten können, definieren Sie zwei Makros. Diese übernehmen das Styling und setzen den Zipfel automatisch links bzw. rechts.

\begin{lstlisting}[style=latex]
\newcommand{\lehrerblase}[2][]{%
  \tikz[baseline=(B.base)]{
    \node[
      rectangle callout,
      rounded corners=6pt,
      callout relative pointer={(-13mm,-7mm)},
      draw, thick,
      fill=blue!6,
      align=left,
      inner sep=3.5mm,
      #1
    ] (B) {#2};
  }%
}

\newcommand{\schuelerblase}[2][]{%
  \tikz[baseline=(B.base)]{
    \node[
      rectangle callout,
      rounded corners=6pt,
      callout relative pointer={(13mm,-7mm)},
      draw, thick,
      fill=green!6,
      align=left,
      inner sep=3.5mm,
      #1
    ] (B) {#2};
  }%
}
\end{lstlisting}

\noindent Verwendung im Dokument:

\begin{lstlisting}[style=latex]
\lehrerblase[text width=0.82\linewidth]{\textbf{Impuls:} Was passiert bei einer Verschiebung um 3?}

\vspace{5mm}

\hfill\schuelerblase[text width=0.62\linewidth]{Man ersetzt jeden Buchstaben durch den dritten danach.}
\end{lstlisting}

\begin{ausgabe}
\lehrerblase[text width=0.82\linewidth]{\textbf{Impuls:} Was passiert bei einer Verschiebung um 3?}

\vspace{5mm}

\hfill\schuelerblase[text width=0.62\linewidth]{Man ersetzt jeden Buchstaben durch den dritten danach.}
\end{ausgabe}

\noindent\textbf{Hinweis zur Praxis:} Callout-Nodes eignen sich besonders gut für kurze bis mittlere Notizen. Wenn Sie sehr lange Texte in einer „Blase“ über Seiten umbrechen möchten, ist eine alternative Konstruktion (z.\,B. mit \texttt{tcolorbox} + eigener Zipfel-Zeichnung) sinnvoll.

\end{document}
